\section{Riesgos y agenda de investigación: qué hacer a continuación en el entrenamiento abierto}

\subsection{Las tres brechas de investigación más críticas en la actualidad}
\begin{table}[H]
\centering
\renewcommand{\arraystretch}{1.2}
\begin{tabular}{p{3.0cm}p{4.0cm}p{4.2cm}}
\toprule
\textbf{Brecha} & \textbf{Problema actual} & \textbf{Dirección siguiente} \\
\midrule
Ecosistema de verificadores & La cobertura de tareas verificables sigue siendo limitada y la transferencia entre dominios es débil & Construir bibliotecas de verifiers componibles y estándares de estratificación de tareas \\
Gobernanza de datos & La calidad de los datos sintéticos es inestable y la contaminación es difícil de auditar de manera unificada & Unificar protocolos de descontaminación y cadenas de herramientas de auditoría de datos de código abierto \\
Control de costos & Los beneficios y la predicción de costos del escalamiento en tiempo de prueba son inestables & Establecer una estrategia de programación conjunta de ``calidad-latencia-costo'' \\
\bottomrule
\end{tabular}
\caption{Tres agendas de investigación aplicables derivadas del contenido de la clase}
\end{table}

\begin{knowledgebox}{Implicaciones directas para los equipos de ingeniería}
Documentar y versionar la ``receta de entrenamiento'', gestionándola con el mismo rigor que el código. Cada cambio de capacidad debe poder rastrearse hasta una modificación concreta en los datos, la función objetivo, el optimizador o la estrategia de inferencia.
\end{knowledgebox}

\subsection{Dos recomendaciones para investigadores y equipos de producto}
\begin{importantbox}{Recomendación 1: priorizar la construcción de una base experimental reproducible}
Aunque a corto plazo no se persiga la puntuación más alta, primero hay que dejar funcionando el entrenamiento, la evaluación y las pruebas de regresión. Sin una base reproducible, cualquier mejora puede resultar insostenible.
\end{importantbox}

\begin{importantbox}{Recomendación 2: tratar el presupuesto de inferencia como un parámetro de producto}
El escalamiento en tiempo de inferencia no es un accesorio algorítmico, sino una estrategia de producto. El presupuesto debe asignarse por niveles según el valor de la solicitud, en lugar de usar la misma estrategia de decodificación para todas las solicitudes.
\end{importantbox}

\begin{warningbox}{La apertura no significa estándares bajos}
La ruta del entrenamiento abierto también requiere un cumplimiento normativo de datos riguroso, descontaminación de la evaluación y revisión de seguridad. De lo contrario, la apertura terminará replicando ruido en lugar de replicar progreso.
\end{warningbox}

\subsection{Resumen de la sección}
El factor decisivo de la siguiente etapa del entrenamiento abierto no es un único algoritmo nuevo, sino esta infraestructura triple: el ecosistema de verificadores, la gobernanza de datos y la programación del presupuesto de inferencia.

